\subsection{Ebene durch drei Punkte}
\label{subsec:analytic_geometry_plane_through_three_points}

Drei nicht auf einer Geraden liegende Punkte bestimmen eine Ebene eindeutig.
Verbindungsvektoren zwischen den Punkten liefern unmittelbar zwei geeignete
Richtungen.

\begin{KeyPointBox}{Ebene durch drei Punkte}{analytic_geometry_plane_from_three_points}
  Liegen die Punkte $A$, $B$ und $C$ nicht auf einer gemeinsamen Geraden, so
  bestimmen sie eindeutig eine Ebene.

  \begin{minipage}[t]{0.51\linewidth}
    \vspace{0pt}
    \textbf{Parameterdarstellung in drei Schritten:}
    \begin{enumerate}[leftmargin=1.6em,itemsep=0.25em]
      \item Aufpunkt wählen, z.\,B. $A$: $\overrightarrow{OA}$.
      \item Richtungen wählen:
        $\vec b\parallel\overrightarrow{AB}$ und
        $\vec c\parallel\overrightarrow{AC}$.
      \item Parameterform:
        $E:\ \vec x=\overrightarrow{OA}+\lambda\vec b+\mu\vec c$,
        $\lambda,\mu\in\R$.
    \end{enumerate}
  \end{minipage}\hfill
  \begin{minipage}[t]{0.46\linewidth}
    \vspace{0pt}\centering
    \CachedTikzFigure{content/resources/chapters/analytic_geometry/tikz/plane_through_three_points.tex}
  \end{minipage}

\end{KeyPointBox}

\begin{ExampleBox}{Ebene durch drei Punkte bestimmen}{analytic_geometry_plane_from_points_example}
  Wir suchen die Ebene $E$ in Parameterform, die durch die folgenden drei
  Punkte verläuft:
  \[
    A(6,3,-5),\qquad B(2,-1,3),\qquad C(1,2,4).
  \]
  \IfClozeTF{%
    \begin{ClozeArea}[34]{20}
    \end{ClozeArea}%
  }{%
    Wir wählen $A$ als Aufpunkt und erhalten zunächst den Stützvektor
    \[
      \overrightarrow{OA}=\begin{pmatrix}6\\3\\-5\end{pmatrix}.
    \]
    Die Richtungsvektoren wählen wir in Richtung der Verbindungsvektoren
    $\overrightarrow{AB}$ und $\overrightarrow{AC}$:
    \begin{ClozeArea}[34]{10}
      \[
        \begin{alignedat}{3}
          \overrightarrow{AB} = \overrightarrow{OB} - \overrightarrow{OA}
            &=\begin{pmatrix}2\\-1\\3\end{pmatrix}
              -\begin{pmatrix}6\\3\\-5\end{pmatrix}
              =\begin{pmatrix}-4\\-4\\8\end{pmatrix}
              =-4\begin{pmatrix}1\\1\\-2\end{pmatrix}
            &\qquad&\text{wähle}\quad
            &\vec b&=\begin{pmatrix}1\\1\\-2\end{pmatrix},
              \quad \text{da } \vec b\parallel\overrightarrow{AB},\\[1em]
          \overrightarrow{AC} = \overrightarrow{OC} - \overrightarrow{OA}
            &=\begin{pmatrix}1\\2\\4\end{pmatrix}
              -\begin{pmatrix}6\\3\\-5\end{pmatrix}
              =\begin{pmatrix}-5\\-1\\9\end{pmatrix}
            &\qquad&\text{wähle}\quad
            &\vec c&=\begin{pmatrix}-5\\-1\\9\end{pmatrix},
              \quad \text{da } \vec c\parallel\overrightarrow{AC}.
        \end{alignedat}
      \]
    \end{ClozeArea}
    Damit $\vec b$ und $\vec c$ tatsächlich eine Ebene aufspannen, müssen sie
    linear unabhängig sein. Für 2 Vektoren ist lineare
    Abhängigkeit gleichbedeutend mit einer kollinearen Ausrichtung, d.h.  \[
      \vec c=r\vec b
      \quad\Longrightarrow\quad
      \begin{cases}
        -5=r,\\
        -1=r
      \end{cases}
      \quad
      \mathrel{\textcolor{red!75!black}{%
        \raisebox{-0.35em}{\fontsize{27}{27}\selectfont $\lightning$}}}
      \quad\Longrightarrow\quad
      \vec b\nparallel\vec c
      \quad\Longrightarrow\quad
      \vec b,\vec c\ \text{sind linear unabhängig}.
    \]
    Eine Parameterdarstellung der gesuchten Ebene lautet daher
    \[
      E:\quad
      \vec x
      =\begin{pmatrix}6\\3\\-5\end{pmatrix}
      +\lambda\begin{pmatrix}1\\1\\-2\end{pmatrix}
      +\mu\begin{pmatrix}-5\\-1\\9\end{pmatrix},
      \qquad \lambda,\mu\in\R.
    \]
  }%
\end{ExampleBox}
